\section{实验}
\label{sec:experiments}
我们在仿真和真机上比较 \ours 与反应式、动态及世界—动作策略（第~\ref{sec:main_results}~节与第~\ref{sec:real}~节），然后对目标历史与未来监督做消融（第~\ref{sec:principles}~节）。

\subsection{设置}
\label{sec:setup}

\paragraph{\bench。}
我们在 MuJoCo/robosuite \cite{mujoco,robosuite} 中构建了五个任务族共 14 个任务：抓取运动目标（简谐、传送带、平移、旋转与遮挡运动）、放入移动的容器、让棒沿旋转的导线穿过、捕获滚动或滑动的物体，以及撞击后的抓取（附录~\appref{app:benchmark}，图~\appref{fig:benchmark_tasks}）。我们在原任务运动范围内扫描附加延迟（0--200\,ms）。使用仿真特权状态的专家为每个任务提供 \cfgDemos{} 条示范。我们在开发场景上选择配置，并在 1{,}400 个留出场景（每个任务 100 个，新种子）上报告 \bench 主要比较与最终模型消融。最终模型的消融取三个训练种子的平均；筛选研究的种子与划分分别见附录~\appref{app:design}。每个基线训练一次。我们给出对留出场景自助法得到的配对 95\% 区间。

\paragraph{协议。}
观测、场景运动与执行共用一个控制周期为 50\,ms 的物理时钟。所有方法在同一执行器下运行，方法计算期间物理世界按其在一块 \cfgGPU{} 上测得的延迟继续推进。一条指令执行完毕且上一次推理结束后，策略再次被查询；新动作块到达之前，执行器继续执行上一个动作块。+150\,ms 设置为每个方法额外加入 150\,ms 延迟，暂停设置在推理期间停止物理仿真（附录~\appref{app:sweeps}）。我们报告按任务等权的成功率（SR）、接触时刻误差 $e_t$、首次接触时的位置误差 $e_p$ 以及 p95 出动作时间。

\paragraph{基线。}
我们比较反应式策略（ACT \cite{act}、Diffusion Policy \cite{diffusionpolicy} 与 $\pi_{0.5}$ \cite{pi05}）、动态 VLA（DynamicVLA \cite{dynamicvla} 与 PUMA \cite{domino}）、图像或隐空间 WAM（DreamVLA \cite{dreamvla} 与 OpenWAM-$\alpha$ \cite{openwam_alpha}），以及一个动态 WAM（DynamicWAM \cite{dynamicwam}）。每个基线使用其官方代码、输入与训练配方，在相同示范上训练；若其配方使用腕部相机，则提供渲染的腕部图像（附录~\appref{app:evaluation}）。\ours 只使用主视角，但还使用深度和目标位姿。为了控制目标运动信息的可用性，$\pi_{0.5}$ 与 DynamicWAM 还额外接收我们的 SE(3) 历史 token（+SE(3)）；它们的深度输入仍与本文方法不同。\emph{匀速输入}对照是去掉预测损失、改为接收目标位姿匀速外推作为输入 token 的 \ours。

% Transcribed from results/realtime_c/ (INFER0): ACT/DP native/+150 ms/Paused SR and 5090 p95. Other results pending.
% Highlight reported values per column with \bestcell/\secondcell (SR: highest; p95: lowest); exclude pending cells.
\begin{table}[t]
\centering\footnotesize
\caption{\tablelang{\textbf{Main results.} Task-balanced SR on 1{,}400 held-out \bench scenes: native latency, +150\,ms, and paused physics. Three training seeds for \ours, one per baseline. p95: time-to-action on \cfgGPU{}. Wrist and history denote extra inputs (Sec.~\ref{sec:setup}). Blue/bold and light blue mark the best and second reported values per column; rankings remain provisional while results are pending.}{\textbf{主要结果。}1{,}400 个 \bench 留出场景的任务均衡 SR：原生延迟、+150\,ms 和暂停物理。\ours 训练三个种子，基线各一个。p95：\cfgGPU{} 上的出动作时间。腕部及历史列表示额外输入（第~\ref{sec:setup}~节）。每列已有结果的最佳值为粗体加蓝底，次佳为浅蓝底；结果待齐，排名为暂时性比较。}}
\label{tab:main}
\setlength{\tabcolsep}{0.8pt}
\begin{tabular}{@{}lccccrrrr@{}}
\toprule
 & & \multicolumn{3}{c}{\tablelang{History}{历史}} & \multicolumn{3}{c}{SR (\%) $\uparrow$} & \\
\cmidrule(lr){3-5}\cmidrule(lr){6-8}
\tablelang{Method}{方法} & \tablelang{Wrist}{腕部} & RGB & \tablelang{Flow}{光流} & SE(3) & \tablelang{Native}{原生} & \makecell{+150\\ms} & \tablelang{Paused}{暂停} & \makecell{p95\\(ms)\\$\downarrow$} \\
\midrule
ACT \cite{act} &  &  &  &  & \bestcell{42.0} & \bestcell{24.1} & \bestcell{40.9} & \bestcell{10.7} \\
DP \cite{diffusionpolicy} &  & \cmark &  &  & \secondcell{0.0} & \secondcell{0.0} & \secondcell{33.1} & \secondcell{1701.9} \\
$\pi_{0.5}$ \cite{pi05} &  &  &  &  & \resultplaceholder & \resultplaceholder & \resultplaceholder & \resultplaceholder \\
$\pi_{0.5}$+SE(3) & \resultplaceholder &  &  & \cmark & \resultplaceholder & \resultplaceholder & \resultplaceholder & \resultplaceholder \\
\midrule
DynamicVLA \cite{dynamicvla} &  & \cmark &  &  & \resultplaceholder & \resultplaceholder & \resultplaceholder & \resultplaceholder \\
PUMA \cite{domino} &  &  & \cmark &  & \resultplaceholder & \resultplaceholder & \resultplaceholder & \resultplaceholder \\
\midrule
DreamVLA \cite{dreamvla} &  & \cmark &  &  & \resultplaceholder & \resultplaceholder & \resultplaceholder & \resultplaceholder \\
OpenWAM-$\alpha$ \cite{openwam_alpha} &  &  &  &  & \resultplaceholder & \resultplaceholder & \resultplaceholder & \resultplaceholder \\
\midrule
DynamicWAM \cite{dynamicwam} &  &  & \cmark &  & \resultplaceholder & \resultplaceholder & \resultplaceholder & \resultplaceholder \\
DynamicWAM+SE(3) & \resultplaceholder &  & \cmark & \cmark & \resultplaceholder & \resultplaceholder & \resultplaceholder & \resultplaceholder \\
\midrule
\ours &  &  &  & \cmark & \resultplaceholder & \resultplaceholder & \resultplaceholder & \resultplaceholder \\
\bottomrule
\end{tabular}
\end{table}


\subsection{主要结果}
\label{sec:main_results}
\ours 在 \bench 留出场景上的 SR 为 \fillin{XX.X}\%（表~\ref{tab:main}），比最强基线 \fillin{最强基线名称} 高 \fillin{X.X} 个百分点，比接收相同 SE(3) 历史的最强基线 \fillin{名称} 高 \fillin{X.X} 个百分点。优势最大的是 \fillin{旋转目标以及快速的传送带和平移目标}，即领先时长内目标移动最远的任务。它的 p95 出动作时间为 \fillin{XX}\,ms，是图像空间 WAM 的 \fillin{X} 分之一；后者的延迟会叠加到它们需要弥补的领先时长上。加入 150\,ms 延迟后，\ours 保持 \fillin{XX.X}\%，而反应式策略降至 \fillin{XX.X}\%，图像空间 WAM 降至 \fillin{XX.X}\%。推理期间暂停物理仿真时，它与 \fillin{暂停设置下最强基线名称} 的差为 \fillin{$\pm$X.X} 个百分点，从而把速度的贡献与表示的贡献区分开。匀速输入变体落后 \ours \fillin{X.X} 个百分点。我们结合无未来监督对照解释这一差异；后者既没有预测损失，也没有外推输入 token。DOMINO \cite{domino} 上的结果见附录~\appref{app:external}。

\subsection{真机实验}
\label{sec:real}

% Planned results: six-task real-robot study; numeric results and baseline selection pending.
\begin{table}[t]
\centering\footnotesize
\caption{\tablelang{\textbf{Real-robot result template.} Six planned tasks: T1 conveyor grasping, T2 moving-receptacle placement, T3 turntable grasping, T4 moving-car grasping, T5 post-collision interception, and T6 ramp interception. Trial counts and baseline selection await confirmation. Cells reserve success rates (\%); All pools the trials.}{\textbf{真机结果模板。}六个计划任务：T1 传送带抓取、T2 移动容器放置、T3 转盘抓取、T4 移动小车抓取、T5 碰撞后拦截、T6 斜坡拦截。试验数与基线选择待确认。单元格预留成功率（\%）；All 合并全部试验。}}
\label{tab:real}
\setlength{\tabcolsep}{2pt}
\resizebox{\columnwidth}{!}{%
\begin{tabular}{@{}lccccccc@{}}
\toprule
\tablelang{Method}{方法} & T1 & T2 & T3 & T4 & T5 & T6 & All \\
\midrule
\fillin{$\pi_{0.5}$} & \resultplaceholder & \resultplaceholder & \resultplaceholder & \resultplaceholder & \resultplaceholder & \resultplaceholder & \resultplaceholder \\
\fillin{DynamicWAM} & \resultplaceholder & \resultplaceholder & \resultplaceholder & \resultplaceholder & \resultplaceholder & \resultplaceholder & \resultplaceholder \\
\tablelang{\ours, no DROID base}{\ours，无 DROID 基座} & \resultplaceholder & \resultplaceholder & \resultplaceholder & \resultplaceholder & \resultplaceholder & \resultplaceholder & \resultplaceholder \\
\ours & \resultplaceholder & \resultplaceholder & \resultplaceholder & \resultplaceholder & \resultplaceholder & \resultplaceholder & \resultplaceholder \\
\bottomrule
\end{tabular}}
\end{table}


我们在\cfgRobotZh{}上评估六个任务（图~\appref{fig:real_setup}，表~\appref{tab:real_tasks}），比较 \ours、$\pi_{0.5}$、DynamicWAM 与 Diffusion Policy \cite{diffusionpolicy}：传送带抓取（T1）、放入传送带上的容器（T2）、转盘抓取（T3）、抓取运动的玩具小车（T4）、拦截碰撞后的球（T5），以及拦截沿斜坡滚落或滑落的物体（T6）。它们涵盖直线、圆周、平面自由运动、撞击后运动与重力加速运动。每个方法在相同的 100 条遥操作示范上逐任务训练或微调，另用独立的 20 条示范验证，并以共享的试验键在每个任务上运行 20 次试验。\ours 用\fillin{跟踪器}跟踪目标，在一块 \cfgGPU{} 上运行，p95 出动作时间为 \fillin{XX}\,ms。它在全部试验中的成功率为 \fillin{XX.X}\%，\fillin{最强基线}为 \fillin{XX.X}\%（表~\ref{tab:real}）。\fillin{优势最大的任务及主要失败模式}（图~\ref{fig:real_trials}）。

\begin{figure}[t]
\centering
\figplaceholder[1.45in]{\columnwidth}{图 4——真机试验}{%
\textbf{行：}\ours 的一次成功；最强基线在同一试验键上的表现；\ours 的一次失败。
\textbf{列：}查询帧、首个执行时刻 $\tapp$ 与接触时刻。\textbf{叠加：}\textcolor{obscol}{蓝色}为跟踪得到的过去位姿；\textcolor{actcol}{紫色}为指令的夹爪路径。真实照片不做修饰；叠加来自各次试验的日志数据。}
\caption{\textbf{真机试验}：\ours 与\fillin{最强基线}在同一试验键上的表现（上两行），以及 \ours 的一次失败（下）。\fillin{对所选试验的一句话说明。}}
\label{fig:real_trials}
\end{figure}

\subsection{消融：动态策略需要什么？}
\label{sec:principles}
四个核心对照使用最终 24 层模型与配方，每个对照训练三个种子，数据与训练预算相同（表~\ref{tab:design}）。它们检验目标历史、未来监督、场景流监督和专用目标位姿监督。各对照在实时协议下评估 0 与 150\,ms 附加延迟，完整模型的参考结果复用第~\ref{sec:main_results}~节。更广泛的设计选择在 4 层模型上筛选，各组配方见附录~\appref{app:design}。

% Planned results: A4 final 24-layer confirmation, four controls x three training seeds, real-time C at added 0/150 ms. Full-model reference reused from A1; no screening numbers substituted.
\begin{table}[t]
\centering\footnotesize
\caption{\tablelang{\textbf{Core ablations} of the final 24-layer model on held-out \bench scenes (three training seeds). Each row changes one component of
\ours (\sel). Native, +150\,ms: success (\%) with native latency and with 150\,ms of added latency. Bottom: the CV-input alternative, which replaces the forecast losses by constant-velocity extrapolations as input tokens. $e_p$: position
error at first contact (mm). p95: time-to-action (ms).}{\textbf{核心消融}（最终 24 层模型，\bench 留出场景，三个训练种子）。每行只改变 \ours（\sel）的一个组件。
原生、+150\,ms：原生延迟下与附加 150\,ms 延迟下的成功率（\%）；下方为匀速外推输入的替代方案，它以匀速外推作为输入 token 取代预测损失；$e_p$：首次接触时的位置误差（mm）；p95：出动作时间（ms）。}}
\label{tab:design}
\setlength{\tabcolsep}{2.9pt}
\begin{tabular}{@{}lcccc@{}}
\toprule
\tablelang{Variant}{变体} & \tablelang{Native}{原生} & +150\,ms & $e_p$ & p95 \\
\midrule
\ours\sel & \resultplaceholder & \resultplaceholder & \resultplaceholder & \resultplaceholder \\
\tablelang{No target history}{无目标历史} & \resultplaceholder & \resultplaceholder & \resultplaceholder & \resultplaceholder \\
\tablelang{No future supervision}{无未来监督} & \resultplaceholder & \resultplaceholder & \resultplaceholder & \resultplaceholder \\
\tablelang{No scene-flow supervision}{无场景流监督} & \resultplaceholder & \resultplaceholder & \resultplaceholder & \resultplaceholder \\
\tablelang{No dedicated SE(3)}{无专用 SE(3)} & \resultplaceholder & \resultplaceholder & \resultplaceholder & \resultplaceholder \\
\midrule
\tablelang{CV input}{匀速外推输入} & \resultplaceholder & \resultplaceholder & \resultplaceholder & \resultplaceholder \\
\bottomrule
\end{tabular}
\end{table}


\begin{figure*}[t]
\centering
\figplaceholder[1.55in]{\textwidth}{图 5——消融（四个子图）}{%
\textbf{(a) 历史。}完整模型与无目标运动历史在 0 和 150\,ms 附加延迟下的 SR 差值。
\textbf{(b) 未来监督。}完整模型与关闭全部未来损失，按任务族分组。
\textbf{(c) 几何监督。}完整模型、无场景流和无专用 SE(3)，按任务族分组。
\textbf{(d) 成功率与延迟。}完整模型与四个对照在 0 和 150\,ms 附加延迟下的 p95 出动作时间与 SR。
各面板复用相同留出轨迹与三个训练种子，显示 95\% 配对自助法区间。}
\caption{\textbf{最终模型规模上的核心消融。}(a) 目标运动历史。(b) 未来监督。(c) 两项几何监督。(d) 实时协议下的成功率与延迟。\fillin{四个对照完成后概括实测效应。}}
\label{fig:analysis}
\end{figure*}

\paragraph{记住什么。}
无历史对照保留当前图像、当前目标位姿和八个机器人状态，移除过去的目标位姿及其速度、加速度特征。它相对完整模型的 SR 差值在 0\,ms 时为 \fillin{带符号 X.X} 个百分点，在 150\,ms 附加延迟时为 \fillin{带符号 X.X} 个百分点（图~\ref{fig:analysis}a）。视觉历史替代方案属于 4 层筛选，不是最终规模的对照。位姿噪声鲁棒性单独用冻结权重评测（附录~\appref{app:design_remember}）。

\paragraph{预测什么。}
无未来对照关闭 SE(3)、场景流及可见性、未来深度、DINO 与动态 RGB 损失，保留当前目标 RGB 重建，用于检验未来监督的整体作用（图~\ref{fig:analysis}b）。另外两个对照保留全部外观目标：一个只移除场景流与可见性监督，另一个只移除专用 SE(3) 监督。后者仍用同一 \cfgContextGrid{} 流覆盖目标、机器人与周围环境，场景各部分均使用相同的粗粒度表示。前者检验场景流监督的作用，不能单独归因于环境信息。两者的带符号 SR 差值在 0\,ms 时为 \fillin{X.X} 和 \fillin{X.X} 个百分点，在 150\,ms 时为 \fillin{X.X} 和 \fillin{X.X} 个百分点；图~\ref{fig:analysis}c 按任务族分解这些效应。

\begin{figure}[t]
\centering
\figplaceholder[1.35in]{\columnwidth}{图 6——执行时刻的预测}{%
\textbf{行：}一次滚动捕获（预测误差为中位数）与一次失败的撞击后抓取。\textbf{列：}$\tau_t$ 的查询帧、$\tapp$ 的帧，以及基座坐标系中的 3D 小图。
\textbf{叠加：}\textcolor{obscol}{蓝色}为过去位姿，\textcolor{predcol}{橙色}为训练期预测头给出的预测三轴与粗粒度流箭头，\textcolor{gtcol}{灰色}为同一时刻的真值，\textcolor{actcol}{紫色}为指令路径。完整示例见图~\appref{fig:qualitative}。}
\caption{\textbf{执行时刻的预测}（\bench 留出回合，读自训练期预测头）。上：SE(3) 预测跟随滚动目标，而外推冲过了头。下：撞击之前，回归得到的预测对可能的结果取平均，抓取落空。}
\label{fig:forecasts}
\end{figure}

\paragraph{如何使用预测。}
最终模型只把预测用于训练监督。4 层耦合研究的三个种子中，把预测送给动作没有带来稳定增益（附录~\appref{app:design_act}）；该结果不能证明最终规模下的等效性。最终模型开启未使用的预测头时，只检查固定输入与噪声下的动作一致性并测量额外代价。实时协议下，额外延迟可能改变闭环轨迹。图~\ref{fig:analysis}d 报告四个训练对照的实测成功率与延迟，不包含最终规模的 R5 比较。
